\documentclass[11pt]{article}
\usepackage[spanish]{babel}
\usepackage[utf8]{inputenc}
\usepackage[T1]{fontenc}
\usepackage{lmodern}
\usepackage{geometry}
\usepackage{amsmath,amssymb,amsfonts}
\usepackage{enumitem}
\usepackage{hyperref}
\usepackage{xcolor}
\usepackage{titlesec}
\geometry{margin=2.6cm}
\hypersetup{colorlinks=true,linkcolor=blue,citecolor=blue,urlcolor=blue}
\setlength{\parskip}{0.6em}
\setlength{\parindent}{0pt}
\titleformat{\section}{\large\bfseries}{\thesection.}{0.6em}{}
\titleformat{\subsection}{\normalsize\bfseries}{\thesubsection.}{0.6em}{}

\title{\vspace{-1cm}
Programa de investigación conjunta (informe rápido)\\[0.3em]
\large Random Matrix Theory y Genómica Computacional\\
\normalsize Propuesta de colaboración con el Dr. Enrique Hernández-Lemus}
\author{}
\date{Septiembre de 2026}

\begin{document}
\maketitle
\thispagestyle{empty}

\section*{Resumen ejecutivo}
Se propone un programa de investigación conjunta que use \emph{Random Matrix Theory} (RMT) --- estática (ensambles, umbrales espectrales) y dinámica (movimiento browniano de Dyson, repulsión de niveles) --- junto con análisis multiescala por \emph{wavelets}, como marco matemático para problemas actuales de genómica computacional y oncología de sistemas, con foco en cáncer de mama. La idea central es modelar la frontera entre \emph{señal biológica estructurada} y \emph{ruido de alta dimensión} en matrices derivadas de datos multi-ómicos, redes de coexpresión a gran escala, matrices de contacto cromatínico tipo Hi-C y datos de célula única, y validar esa frontera con protocolos numéricos explícitos y reproducibles sobre datos públicos reales. La propuesta está alineada con la trayectoria reciente del Dr. Enrique Hernández-Lemus en integración multi-ómica en cáncer, redes regulatorias multicapa, genómica 3D y redes específicas por paciente.

\section{Justificación científica}
El Dr. Enrique Hernández-Lemus dirige el Laboratorio de Biología Computacional de Sistemas y Genómica Integrativa del Departamento de Genómica Computacional del INMEGEN, con trabajo reciente en integración de datos multi-ómicos en cáncer \cite{multiomic2022}, redes regulatorias multicapa \cite{multilayer2021}, estructura de \emph{k-cores} en redes de cáncer de mama \cite{kcore2021}, conformación cromosómica en cáncer de mama triple negativo \cite{tnbc2025} y redes de coexpresión específicas por paciente \cite{luad2025}. Estas líneas sugieren un contexto ideal para introducir herramientas espectrales rigurosas --- estáticas y dinámicas --- de RMT sobre matrices de covarianza, correlación, afinidad y contacto cromatínico \cite{inmegen, multiomic2024, tnbc2025, luad2025}.

Desde el punto de vista metodológico, RMT ofrece una ventaja estratégica: permite construir \emph{ensambles nulos} y umbrales espectrales no ad hoc para distinguir estructura biológica real frente a artefactos debidos a alta dimensionalidad, escasez muestral, heterocedasticidad y sparsidad. El método fundacional de Luo et al.\ \cite{luo2007}, ya validado a gran escala \cite{gao2012} y extendido a redes de interacción génica en cáncer \cite{cancer2018}, usa precisamente la transición de la distribución de espaciamientos entre vecinos más cercanos (NNSD), de Poisson a la estadística de Wigner-Dyson del ensamble ortogonal gaussiano (GOE), para fijar un umbral de correlación no arbitrario en redes de coexpresión --- exactamente el tipo de red a gran escala (\emph{big scale networks}) que el grupo del Dr.\ Hernández-Lemus ya construye. Esta lógica ya ha mostrado utilidad en análisis de célula única, donde existen trabajos de denoising guiado por RMT y métodos espectrales para extraer señal de matrices extremadamente ruidosas \cite{aparicio2020, pnas2021}. Asimismo, en Hi-C reciente se ha señalado que ciertos esquemas de suavizado por random walks deben tratarse con cautela, lo cual abre espacio para null models espectrales mejor fundamentados \cite{hic2025}.

\section{Hipótesis rectora}
La hipótesis del programa es la siguiente:
\begin{quote}
\emph{La estructura biológica relevante puede detectarse y cuantificarse como desviación espectral robusta --- estática y dinámica --- respecto de un ensamble nulo de alta dimensión, y esa desviación es numéricamente validable sobre cohortes independientes.}
\end{quote}

Dicho de otro modo, autovalores salientes, autovectores localizados, estadísticos lineales de eigenvalores, rupturas universales del \emph{bulk}, y patrones de repulsión/casi-colisión entre autovalores bajo perturbación o evolución temporal, pueden funcionar como marcadores matemáticamente controlados de modularidad funcional, heterogeneidad tumoral, reorganización cromatínica y señal paciente-específica --- siempre que la desviación sobreviva a un protocolo explícito de validación estadística.

\section{Objetivo general}
Desarrollar una agenda conjunta de investigación que integre Random Matrix Theory (estática y dinámica), análisis multiescala por wavelets, teoría espectral de redes, y aprendizaje automático, con datos reales de genómica computacional, con énfasis en cáncer de mama, multi-ómica, genómica 3D, célula única y redes a gran escala, y con un protocolo de validación numérica explícito en cada línea.

\section{Objetivos específicos}
\begin{enumerate}[label=\arabic*.]
    \item Construir modelos nulos espectrales para matrices de expresión, coexpresión y covarianza multi-ómica, a la escala de redes completas del transcriptoma (\(p\sim 10^4\)--\(10^5\)).
    \item Diseñar criterios RMT --- estáticos (NNSD, IPR) y dinámicos (Dyson, repulsión de niveles) --- para detección de módulos biológicos, subtipos tumorales y biomarcadores robustos.
    \item Extender el marco a matrices de contacto cromatínico (Hi-C/scHi-C) y redes multicapa expresión--conformación, incorporando descomposición \emph{wavelet} multiescala como paso de pre-procesamiento y de extracción de características.
    \item Desarrollar variantes para redes específicas por paciente y para datos de célula única de alta sparsidad, y acoplarlas a clasificadores de aprendizaje automático cuyo propio proceso de entrenamiento se analice con las mismas herramientas espectrales (reconocimiento de subtipo/patrón).
    \item Fijar, para cada línea, un protocolo de validación (remuestreo, cohortes independientes, control de falsos positivos) y traducir estos desarrollos a \emph{pipelines} computacionales reproducibles y a artículos de metodología aplicada.
\end{enumerate}

\section{Líneas de trabajo propuestas}

\subsection{Línea 1. Denoising espectral y detección de señal en datos multi-ómicos}
Sea $X \in \mathbb{R}^{p\times n}$ una matriz de datos ómicos centrada y escalada, con $p$ variables y $n$ muestras. La matriz
\[
S=\frac{1}{n}XX^{\top}
\]
puede interpretarse como matriz de covarianza o correlación empírica. El primer objetivo es estudiar el espectro de $S$ frente a ensambles nulos apropiados (Wishart, spiked models, variantes heterocedásticas o con ruido correlacionado), para separar:
\begin{itemize}
    \item \textbf{bulk espectral:} parte compatible con ruido de alta dimensión, delimitada por la ley de Marchenko--Pastur;
    \item \textbf{outliers:} componentes compatibles con señal biológica organizada;
    \item \textbf{localización de autovectores:} firmas de módulos génicos o ejes tumorales, cuantificadas por el \emph{inverse participation ratio} (IPR).
\end{itemize}

\textbf{Meta concreta.} Desarrollar un método de filtrado espectral para matrices multi-ómicas y comparar su desempeño con PCA estándar, sparse PCA y enfoques de redes usados en la literatura reciente de integración multi-ómica \cite{multiomic2024, multiomic2022}.

\subsection{Línea 2. RMT para redes de coexpresión a gran escala y biomarcadores específicos por paciente}
Los trabajos recientes sobre redes específicas por paciente muestran que la estructura regulatoria individualizada puede revelar subtipos y biomarcadores no visibles en redes promedio \cite{luad2025}, y que la estructura de \emph{k-cores} organiza jerárquicamente las redes de cáncer de mama \cite{kcore2021}. El método de umbral espectral de Luo et al.\ \cite{luo2007}, ya usado a la escala de redes completas del transcriptoma \cite{gao2012}, es el punto de partida natural para redes de este tamaño. Aquí la contribución conjunta puede ser fuerte: pasar de una red inferida empíricamente a una teoría de \emph{significancia espectral} de esa red, y de ahí a una noción cuantitativa de qué tan robusto es un \emph{k-core} o un módulo bajo la incertidumbre estadística que el propio umbral espectral hace explícita.

\textbf{Problemas propuestos:}
\begin{enumerate}[label=(\alph*)]
    \item estudiar espectros de matrices de adyacencia o Laplacianos de redes de coexpresión a gran escala bajo null models dependientes del grado y de la distribución empírica;
    \item definir estadísticos basados en la separación de autovalores y en la localización de autovectores para clasificar pacientes;
    \item construir indicadores de estabilidad espectral de biomarcadores y de \emph{k-cores} bajo remuestreo y perturbaciones.
\end{enumerate}

\textbf{Meta concreta.} Producir un artículo metodológico sobre \emph{RMT for patient-specific regulatory networks}, con aplicación a cáncer de mama.

\subsection{Línea 3. Modelos espectrales para Hi-C y genómica 3D}
La integración entre conformación cromosómica y expresión génica ya es una línea activa en el grupo del Dr.\ Hernández-Lemus, particularmente en cáncer de mama triple negativo \cite{tnbc2025}. Aquí RMT puede aportar una base matemática novedosa para matrices de contacto, kernels de proximidad y operadores derivados de Hi-C.

\textbf{Preguntas clave:}
\begin{itemize}
    \item ¿Qué parte del espectro de una matriz de contacto cromatínico corresponde a arquitectura robusta y cuál a ruido experimental o a profundidad de secuenciación?
    \item ¿Cómo definir ensambles nulos que respeten sesgos marginales, dependencia con la distancia genómica y sparsidad?
    \item ¿Qué estadísticos espectrales detectan reordenamientos estructurales en tumores frente a tejido normal?
\end{itemize}

\textbf{Meta concreta.} Desarrollar una teoría preliminar de \emph{spectral null models for chromatin contact matrices} y aplicarla a datos TNBC. Esto conectaría de manera directa con el trabajo reciente del grupo y le daría una capa matemática diferenciadora.

\subsection{Línea 4. RMT en célula única y multimodalidad}
El análisis de célula única y topología de datos ya aparece como área de expansión natural del grupo \cite{csbig, pnas2021}. Los datos scRNA-seq y multiome presentan precisamente el régimen donde RMT es más valiosa: gran dimensión, matrices esparsas, ruido técnico elevado y señal de baja jerarquía.

\textbf{Metas concretas:}
\begin{enumerate}[label=(\alph*)]
    \item adaptar técnicas de denoising RMT a scRNA-seq y scATAC/scRNA integrados;
    \item usar umbrales Tracy--Widom o pruebas tipo Marchenko--Pastur para selección de componentes informativos;
    \item estudiar el acoplamiento entre RMT y métodos topológicos, usando los objetos espectralmente depurados como entrada para TDA o clustering jerárquico.
\end{enumerate}

\subsection{Línea 5. Análisis multiescala por wavelets}
Las redes de coexpresión y las matrices de contacto Hi-C tienen estructura a múltiples escalas simultáneas: módulos génicos locales, dominios topológicamente asociados (TADs) de tamaño intermedio, y compartimentos A/B de gran escala. Un umbral espectral único, aplicado a la matriz completa, puede ser demasiado grueso para separar señal a todas estas escalas a la vez. König et al.\ \cite{konig2006} ya mostraron, sobre la red metabólica de \emph{E.\ coli}, que una transformada wavelet aplicada a los perfiles de expresión --- y a las matrices de adyacencia de los clusters resultantes --- extrae patrones funcionales que un análisis de escala única no separa limpiamente.

\textbf{Propuesta concreta.} Antes de aplicar el umbral RMT de la Línea 2, descomponer la matriz de correlación (o la señal de contacto Hi-C a lo largo de la distancia genómica) en una jerarquía de escalas mediante una transformada wavelet 2D discreta, y aplicar el criterio NNSD/Marchenko--Pastur \emph{por escala}: la hipótesis de trabajo es que el nivel de escala en el que primero aparece una transición Poisson\,$\to$\,Wigner-Dyson identifica la escala característica de la organización modular real (génica, TAD, o compartimento), en vez de imponerla a priori.

\textbf{Meta concreta.} Un método reproducible de \emph{umbral espectral multiescala}, comparado contra el umbral de escala única de Luo et al.\ \cite{luo2007} sobre las mismas cohortes.

\subsection{Línea 6. Dinámica de Dyson, repulsión de niveles, y reconocimiento}
Todo lo anterior es \emph{estático}: se calcula el espectro de una matriz fija. Dyson \cite{dyson1962} mostró que, si las entradas de una matriz aleatoria evolucionan por difusión browniana, sus autovalores evolucionan como un gas de Coulomb unidimensional: cada par de autovalores se repele con fuerza $1/(\lambda_i-\lambda_j)$, de modo que dos autovalores casi nunca colisionan (\emph{level repulsion}, la misma repulsión que produce la estadística de Wigner-Dyson en el caso estático). Esto sugiere dos extensiones concretas, ninguna de las dos presente en las Líneas 1--4:

\begin{enumerate}[label=(\alph*)]
    \item \textbf{Dinámica sobre cohortes/tiempo.} Cuando la matriz de coexpresión o de contacto cromatínico se recalcula a lo largo de una progresión (estadios tumorales, pseudotiempo en célula única, o remuestreo secuencial de una cohorte), la trayectoria de cada autovalor puede tratarse como una realización de un proceso tipo Dyson. Una \emph{casi-colisión} entre dos autovalores --- un mínimo local pronunciado de $|\lambda_i-\lambda_j|$ --- es candidata natural a marcar una transición estructural (p.\,ej.\ una reorganización modular o un cambio de subtipo), de forma análoga a como el mismo fenómeno se usa para detectar transiciones espectrales en otros sistemas físicos.
    \item \textbf{Dinámica del propio reconocedor.} Aarts, Hajizadeh, Lucini y Park \cite{aarts2024} muestran, para redes neuronales entrenadas por descenso de gradiente estocástico, que el espectro de las matrices de pesos evoluciona --- desde una distribución inicial tipo Marchenko--Pastur hacia un gas de Coulomb con potencial dependiente del modelo --- exactamente como movimiento browniano de Dyson, con el nivel de estocasticidad fijado por la razón entre tasa de aprendizaje y tamaño de minibatch. Esto da un lenguaje espectral \emph{común} para los datos (Líneas 1--3) y para el clasificador de reconocimiento de subtipo/patrón que se entrena sobre ellos (Línea 4): en ambos casos, la pregunta de si la estructura aprendida es señal genuina o artefacto de optimización se plantea, literalmente, en los mismos términos de repulsión de autovalores.
\end{enumerate}

\textbf{Meta concreta.} Una nota metodológica --- \emph{Dyson dynamics for spectral transitions in cancer genomics} --- que formalice (a) sobre redes de coexpresión a través de estadios TNBC, y (b) sobre la dinámica de entrenamiento de un clasificador de subtipo entrenado sobre las mismas redes, verificando si ambas trayectorias muestran el mismo régimen de repulsión.

\section{Validación y numérica}
\label{sec:validacion}
Cada línea anterior produce un umbral, un estadístico, o una trayectoria espectral; ninguno se reporta sin un protocolo de validación explícito:
\begin{enumerate}[label=\arabic*.]
    \item \textbf{Ensamble nulo correcto.} Comparación contra Wishart puro \emph{y} contra nulos que preservan momentos empíricos (grado, varianza marginal, sparsidad) --- un umbral que sólo supera al Wishart puro no basta.
    \item \textbf{Remuestreo.} Bootstrap y submuestreo de pacientes para reportar un intervalo, no un único valor, para cada autovalor saliente y cada estadístico de localización (IPR).
    \item \textbf{Cohortes independientes.} Todo umbral o módulo detectado en una cohorte (p.\,ej.\ TCGA-BRCA) se reevalúa, sin reajustar parámetros, en al menos una cohorte independiente (p.\,ej.\ METABRIC o un GSE específico de TNBC); ver \S\ref{sec:datos}.
    \item \textbf{Control de falsos positivos.} Corrección explícita por comparaciones múltiples al reportar módulos, biomarcadores o transiciones de Dyson candidatas.
    \item \textbf{Numérica de referencia.} Un script de referencia (Python/NumPy/SciPy) que calcule, sobre datos sintéticos con espectro conocido, el umbral NNSD, la ley de Marchenko--Pastur teórica y el IPR, como prueba de correctitud antes de aplicar el método a datos reales; ver \texttt{scripts/} en el repositorio de este programa.
\end{enumerate}

\section{Núcleo matemático del programa}
El programa puede organizarse alrededor de seis bloques teóricos:
\begin{enumerate}[label=\arabic*.]
    \item ensambles de Wishart y leyes tipo Marchenko--Pastur;
    \item modelos \emph{spiked} y transición BBP para señal de bajo rango;
    \item estadísticos lineales de eigenvalores y fluctuaciones tipo Tracy--Widom;
    \item localización/delocalización de autovectores y robustez frente a ruido estructurado;
    \item matrices aleatorias sobre grafos y operadores espectrales en redes multicapa a gran escala;
    \item movimiento browniano de Dyson, repulsión de niveles, y análisis multiescala por wavelets como pre-procesamiento espectral.
\end{enumerate}

Esto permite una división natural del trabajo: una parte más teórica (RMT, límites, umbrales, null models, dinámica de Dyson) y otra parte aplicada (pipelines, validación sobre cohortes, interpretación biológica, contraste con estado del arte).

\section{Datos}
\label{sec:datos}
Conjuntos de datos públicos, ya identificados, suficientes para arrancar las Líneas 1--2 y 5--6 sin esperar datos propios del grupo; ver \texttt{data/README.md} en el repositorio de este programa para instrucciones de descarga y notas de preprocesamiento.
\begin{itemize}
    \item \textbf{GSE171958} (GEO) --- metilación de DNA y RNA-seq para TNBC, con proteómica pareada en ProteomeXchange PXD025238; líneas celulares MCF10A, MDA-MB-231, HCC1937 \cite{tnbcmultiomics2021}. Punto de entrada natural para la Línea 1 (denoising multi-ómico) y la Línea 5 (wavelets).
    \item \textbf{GSE176078} (GEO) --- scRNA-seq de tumores de 10 pacientes con TNBC \cite{tnbcsc2021}. Punto de entrada para la Línea 4 (célula única) y la parte (a) de la Línea 6 (dinámica a través de pacientes/pseudotiempo).
    \item \textbf{TCGA-BRCA} --- cohorte pública de referencia (expresión, metilación, mutación, clínico) para cáncer de mama; usada aquí como cohorte primaria de validación cruzada (\S\ref{sec:validacion}, punto 3).
    \item \textbf{METABRIC} --- cohorte pública independiente de TCGA-BRCA, con subtipificación clínica; cohorte de validación cruzada secundaria.
\end{itemize}
Ningún dato clínico o multi-ómico propio del grupo del Dr.\ Hernández-Lemus se asume ni se solicita aquí; estos cuatro conjuntos son de acceso público y bastan para un primer ciclo completo de desarrollo y validación metodológica.

\section{Productos esperados}
\begin{enumerate}[label=\arabic*.]
    \item \textbf{Paper 1:} RMT para filtrado espectral en datos multi-ómicos de cáncer.
    \item \textbf{Paper 2:} significancia espectral en redes específicas por paciente, a gran escala.
    \item \textbf{Paper 3:} null models espectrales para matrices Hi-C y redes multicapa expresión--cromatina, con umbral multiescala por wavelets.
    \item \textbf{Paper 4:} dinámica de Dyson y repulsión de niveles como marcador de transición estructural, con el paralelo dato/clasificador de la Línea 6.
    \item \textbf{Software:} paquete reproducible en \texttt{Python}/\texttt{R} para análisis espectral estático y dinámico de datos genómicos, con protocolo de validación integrado.
    \item \textbf{Seminario conjunto:} puente entre matemáticas, física estadística y genómica computacional.
\end{enumerate}

\section{Ruta tentativa de 12--18 meses}
\subsection*{Fase I (0--4 meses)}
Curación de los datasets de \S\ref{sec:datos}, revisión técnica conjunta, y desarrollo del primer pipeline RMT (Línea 1--2) con el protocolo de validación de \S\ref{sec:validacion} desde el inicio.

\subsection*{Fase II (4--9 meses)}
Aplicación a redes específicas por paciente a gran escala (Línea 2) y elaboración del primer manuscrito metodológico. Paralelamente, diseño de null models multiescala para Hi-C (Líneas 3 y 5).

\subsection*{Fase III (9--14 meses)}
Aplicación a TNBC con datos de conformación cromosómica y expresión (GSE171958, GSE176078); formulación del segundo y tercer artículos. Primera prueba de la dinámica de Dyson (Línea 6a).

\subsection*{Fase IV (14--18 meses)}
Extensión a célula única, prueba de la dinámica de Dyson sobre el propio clasificador (Línea 6b), y consolidación de software, seminarios y siguiente agenda de financiamiento.

\section{Ventaja estratégica de la colaboración}
La colaboración tendría una identidad muy clara: no sería una aplicación superficial de técnicas matemáticas a datos biológicos, sino la construcción de un marco riguroso --- estático y dinámico, de escala única y multiescala --- para \emph{inferencia espectral en genómica de alta dimensión}, con un protocolo de validación explícito en cada paso. El valor agregado para ambas partes es nítido:
\begin{itemize}
    \item para el grupo del Dr.\ Hernández-Lemus, aporta una capa matemática fuerte, original y publicable, directamente sobre sus propias líneas de redes a gran escala, TNBC y multi-ómica;
    \item para la parte matemática, ofrece problemas reales, datasets complejos ya identificados (\S\ref{sec:datos}), y una ruta natural hacia publicaciones interdisciplinarias de alto nivel.
\end{itemize}

\section{Primeros pasos sugeridos}
\begin{enumerate}[label=\arabic*.]
    \item Definir el caso de estudio inicial: TNBC, usando GSE171958 y GSE176078 (\S\ref{sec:datos}) como punto de partida inmediato, sin esperar datos propios.
    \item Correr el script de referencia de \S\ref{sec:validacion} (\texttt{scripts/}) sobre datos sintéticos, y luego sobre una submuestra real, como prueba de concepto conjunta.
    \item Escribir una nota conceptual de 2--3 páginas con una pregunta piloto y la metodología RMT mínima.
    \item Preparar un seminario interno de arranque con lenguaje puente entre genómica y teoría espectral.
\end{enumerate}

\section*{Conclusión}
La colaboración propuesta tiene potencial real para convertirse en una línea estable de investigación. Su fuerza está en que conecta de manera orgánica el trabajo actual del Dr.\ Enrique Hernández-Lemus con una agenda matemática distintiva: RMT --- estática y dinámica --- como teoría de señal, ruido, estructura y robustez en sistemas genómicos complejos, complementada por análisis multiescala y por un protocolo de validación explícito desde el primer día. Bien planteada, esta agenda puede producir resultados conceptuales, metodológicos y computacionales con salida natural hacia bioinformática matemática, biología de sistemas y oncología computacional.

\begin{thebibliography}{99}

\bibitem{inmegen}
INMEGEN, \emph{Dr. Enrique Hernández Lemus}, perfil institucional, Instituto Nacional de Medicina Genómica.
\newblock \url{https://www.inmegen.gob.mx/investigacion/investigadores/curriculum-vitae/?perfil=59}

\bibitem{csbig}
CSB-IG Group, \emph{About Enrique Hernández-Lemus}.
\newblock \url{https://csbig.inmegen.gob.mx/about/}

\bibitem{multiomic2024}
E.~Hernández-Lemus y S.~Ochoa,
\newblock ``Methods for multi-omic data integration in cancer research,''
\newblock \emph{Frontiers in Genetics}, 2024.
\newblock \url{https://www.frontiersin.org/journals/genetics/articles/10.3389/fgene.2024.1425456/full}

\bibitem{multiomic2022}
Grupo Hernández-Lemus,
\newblock ``Functional impact of multi-omic interactions in breast cancer subtypes,''
\newblock \emph{Frontiers in Genetics}, 2022.

\bibitem{multilayer2021}
Grupo Hernández-Lemus,
\newblock ``An Information Theoretical Multilayer Network Approach to Breast Cancer Transcriptional Regulation,''
\newblock \emph{Frontiers in Genetics}, 2021.

\bibitem{kcore2021}
Grupo Hernández-Lemus,
\newblock ``k-core genes underpin structural features of breast cancer,''
\newblock \emph{Scientific Reports}, 2021.

\bibitem{tnbc2025}
H.~Reyes-Gopar, K.~A.~Pérez-Fuentes, M.~L.~Bendall y E.~Hernández-Lemus,
\newblock ``Integration of chromosome conformation and gene expression networks reveals regulatory mechanisms in triple negative breast cancer,''
\newblock \emph{Frontiers in Cell and Developmental Biology}, 2025.
\newblock \url{https://www.frontiersin.org/journals/cell-and-developmental-biology/articles/10.3389/fcell.2025.1597245/full}

\bibitem{luad2025}
P.~López-Sánchez, F.~Ávila-Moreno, E.~Hernández-Lemus, M.~L.~Kuijjer y J.~Espinal-Enríquez,
\newblock ``Patient-specific gene co-expression networks reveal novel subtypes and predictive biomarkers in lung adenocarcinoma,''
\newblock \emph{npj Systems Biology and Applications}, 2025.
\newblock \url{https://www.nature.com/articles/s41540-025-00522-0}

\bibitem{aparicio2020}
L.~Aparicio, K.~Bordyuh, A.~Blumberg y S.~Levit,
\newblock ``A Random Matrix Theory Approach to Denoise Single-Cell Data,''
\newblock \emph{Patterns}, 2020.
\newblock \url{https://pmc.ncbi.nlm.nih.gov/articles/PMC7660363/}

\bibitem{pnas2021}
S.~Levyang et al.,
\newblock ``Revealing lineage-related signals in single-cell gene expression using random matrix theory,''
\newblock \emph{PNAS}, 2021.
\newblock \url{https://www.pnas.org/doi/10.1073/pnas.1913931118}

\bibitem{hic2025}
Y.~Liu et al.,
\newblock ``Can random walking on a Hi-C contact matrix lead to data quality improvement?,''
\newblock 2025.
\newblock \url{https://pmc.ncbi.nlm.nih.gov/articles/PMC12456815/}

\bibitem{luo2007}
F.~Luo, Y.~Yang, J.~Zhong, H.~Gao, L.~Khan, D.~K.~Thompson y J.~Zhou,
\newblock ``Constructing gene co-expression networks and predicting functions of unknown genes by random matrix theory,''
\newblock \emph{BMC Bioinformatics} \textbf{8}, 299 (2007).
\newblock \url{https://doi.org/10.1186/1471-2105-8-299}

\bibitem{gao2012}
Massive-scale gene co-expression network construction and robustness testing using random matrix theory,
\newblock \emph{PLOS ONE}, 2012.
\newblock \url{https://doi.org/10.1371/journal.pone.0055871}

\bibitem{cancer2018}
Random matrix analysis for gene interaction networks in cancer cells,
\newblock \emph{Scientific Reports} \textbf{8} (2018).
\newblock \url{https://doi.org/10.1038/s41598-018-28954-1}

\bibitem{konig2006}
R.~König, G.~Schramm, M.~Oswald et al.,
\newblock ``Discovering functional gene expression patterns in the metabolic network of Escherichia coli with wavelets transforms,''
\newblock \emph{BMC Bioinformatics} \textbf{7}, 119 (2006).
\newblock \url{https://doi.org/10.1186/1471-2105-7-119}

\bibitem{dyson1962}
F.~J.~Dyson,
\newblock ``A Brownian-Motion Model for the Eigenvalues of a Random Matrix,''
\newblock \emph{Journal of Mathematical Physics} \textbf{3}, 1191--1198 (1962).

\bibitem{aarts2024}
G.~Aarts, O.~Hajizadeh, B.~Lucini y C.~Park,
\newblock ``Dyson Brownian motion and random matrix dynamics of weight matrices during learning,''
\newblock preprint, arXiv:2411.13512, 2024. NeurIPS 2024 Workshop on Machine Learning and the Physical Sciences.
\newblock \url{https://arxiv.org/abs/2411.13512}

\bibitem{tnbcmultiomics2021}
Multi-omics data integration reveals correlated regulatory features of triple negative breast cancer,
\newblock \emph{Molecular Omics}, 2021. Data: GEO GSE171958, ProteomeXchange PXD025238.
\newblock \url{https://doi.org/10.1039/D1MO00117E}

\bibitem{tnbcsc2021}
Single-cell RNA-seq of triple negative breast cancer tumors (10 patients),
\newblock GEO accession GSE176078.
\newblock \url{https://www.ncbi.nlm.nih.gov/geo/query/acc.cgi?acc=GSE176078}

\end{thebibliography}

\end{document}
